% Generated draft: requires mathematical and source review.
\documentclass[11pt]{article}
\usepackage[margin=25mm]{geometry}
\usepackage{fontspec}
\setmainfont{Latin Modern Roman}
\usepackage{amsmath,amssymb,mathtools,unicode-math}
\setmathfont{Latin Modern Math}
\usepackage{enumitem,needspace,xurl,hyperref}
\setlength{\parindent}{0pt}
\setlength{\parskip}{5pt}
\setlength{\emergencystretch}{3em}
\newcommand{\N}{\mathbb{N}}
\newcommand{\Z}{\mathbb{Z}}
\newcommand{\Q}{\mathbb{Q}}
\newcommand{\R}{\mathbb{R}}
\newcommand{\C}{\mathbb{C}}
\newcommand{\F}{\mathbb{F}}
\newcommand{\A}{\mathbb{A}}
\newcommand{\PP}{\mathbb{P}}
\newcommand{\E}{\mathbb{E}}
\newcommand{\eps}{\varepsilon}
\newcommand{\dd}{\,\mathrm{d}}
\begin{document}
\section*{UnsolvedMath ID: 30005260}
\textbf{Source problem code:} catalogue entry
\subsection{Definitions and mathematical statement}
The source concerns metric-measure spaces, abbreviated mm-spaces, their isomorphism classes, and certain pseudometrics on those classes. The supplied excerpt does not define the admissible mm-spaces, their isomorphisms, or the pseudometrics. In particular, it does not specify the measure normalization, regularity assumptions, or underlying class of spaces.

For context, a pseudometric on a set \(S\) is a function \(\delta:S\times S\to[0,\infty)\) satisfying, for all \(x,y,z\in S\),
\[
\delta(x,x)=0,\qquad \delta(x,y)=\delta(y,x),\qquad \delta(x,z)\leq\delta(x,y)+\delta(y,z).
\]
It is a metric if it also satisfies \(\delta(x,y)=0\Rightarrow x=y\). The excerpt asks whether its unspecified pseudometrics become metrics on suitable subclasses of isomorphism classes of mm-spaces. This general definition does not identify the particular pseudometrics used in the source.

\textbf{Source question.} Can the curvature restrictions in Theorem 2 of the cited contribution be removed from its metric-measure reconstruction result, while retaining its other hypotheses and its conclusion?

Only a schematic logical formulation is possible from the supplied record. Let \(\mathcal X\) denote the domain of that theorem, let \(H_{\mathrm{other}}(u)\) denote all its hypotheses other than its curvature restrictions, let \(H_{\mathrm{curv}}(u)\) denote those restrictions, and let \(C(u)\) denote its conclusion. Here \(u\) denotes the entire collection of objects quantified over by the theorem. The requested removal would replace an implication of the form
\[
\forall u\in\mathcal X,\qquad
\bigl(H_{\mathrm{other}}(u)\land H_{\mathrm{curv}}(u)\bigr)\Longrightarrow C(u)
\]
by
\[
\forall u\in\mathcal X,\qquad H_{\mathrm{other}}(u)\Longrightarrow C(u).
\]
These symbols are placeholders, not a reconstruction of Theorem 2. Its statement, curvature assumptions, remaining hypotheses, quantifiers, and reconstruction conclusion are absent from the supplied record. Consequently, a complete precise mathematical conjecture cannot be extracted without additional source material.
\subsection{Short English statement}
The question asks whether the curvature assumptions in a cited metric-measure reconstruction theorem can be dropped while keeping its other assumptions and conclusion. The surrounding discussion concerns when certain pseudometrics distinguish isomorphism classes of metric-measure spaces. The supplied excerpt omits Theorem 2 and the relevant definitions, so the formulation remains schematic and requires human mathematical review.
\subsection{Sources}
\begin{itemize}
\item Supplied local UnsolvedMath registry record, permanent ID \(30005260\), source problem code \texttt{catalogue entry}, title: \emph{Removing Curvature Restrictions from Metric Measure Reconstruction}. The record belongs to the Oberwolfach Reports Open Problems collection; no individual proposer is supplied. The proposed year is recorded as 2022.
\item Supplied source citation: \emph{Heat Kernels, Stochastic Processes and Functional Inequalities} (2023), Oberwolfach Reports, pp.\ 2811--2893, DOI \texttt{10.4171/owr/2022/49}. \url{https://doi.org/10.4171/owr/2022/49}. The supplied excerpt identifies report page 2839 and refers to Theorem 2 and reference [6]; neither the theorem nor the bibliographic identity of [6] is supplied.
\item Supplied literature source: \emph{Oberwolfach Report 49}. \url{https://ems.press/content/serial-article-files/46987}.
\item Literature assessment dated exactly \texttt{2026-08-22}: the supplied curation labels the problem open and describes the curvature restrictions in Theorem 2 as possible proof artifacts, with no counterexample or removal theorem established. This is supplied dated curation, not an independently verified literature assessment. The background explicitly requests independent verification of time-sensitive status and cited claims.
\end{itemize}
\end{document}
